\section{Movement required by a specified recovery interface}
\label{sec:movement}

A source cover must satisfy two different requirements: it must contain
enough information to recover the output, and its source positions must
have an affordable joint trajectory.  This section separates those
requirements on a familiar coded-caching example.  The code itself is
known.  The movement bounds concern its specified raw-query interface;
they do not exclude every decoder for that code.

\subsection{A finite bound independent of query order}

Write $\operatorname{dist}_L$ for distance on the circle of circumference
$L>0$.  Fix $N\ge1$ equally spaced centres
$c_i=(i+\tfrac12)L/N$, $0\le i<N$.  A query type is a $k$-element subset
of these centres, where $1\le k\le N$.

\begin{theorem}[Theorem M: movement to cover distinct query types]
\label{thm:movement}
\formal{SourceCovers/QueryMovement.lean}{SourceCovers.QueryMovement.movement_max,
SourceCovers.QueryMovement.movement, SourceCovers.QueryMovement.movement_line,
SourceCovers.QueryMovement.movement_of_charge}
Let $S_1,\ldots,S_J$ be a nonempty sequence of query types, with $F$
distinct types.  For $1\le\ell\le m$, let $y_\ell(j)$ be labelled source
positions on the circle, where $m\ge1$.  Suppose every centre in $S_j$
is within distance $\rho>0$ of some $y_\ell(j)$.  Put
\[
 V=\sum_{\ell=1}^{m}\sum_{j=1}^{J-1}
       \operatorname{dist}_L(y_\ell(j+1),y_\ell(j)),\qquad
 v=V/L,\qquad \lambda=\rho/L,
\]
\[
 s=\min\left\{N,\;m\left(\left\lfloor\frac{5\lambda N}{2}
                                  \right\rfloor+1\right)\right\},
 \qquad B=\binom{s}{k},
\]
with $\binom{s}{k}=0$ for $s<k$.  Then
\begin{equation}
 F\le (1+4v/\lambda)B.
 \label{eq:movement-finite}
\end{equation}
In particular $B>0$, and
\begin{equation}
 v\ge\frac{\lambda}{4}\left(\frac{F}{B}-1\right).
 \label{eq:movement-lower}
\end{equation}
More sharply, put
\[
 C=\sum_{j=1}^{J-1}\max_{1\le\ell\le m}
       \operatorname{dist}_L(y_\ell(j+1),y_\ell(j)).
\]
Then $C\le V$ and
\[
 F\le\left(\left\lfloor\frac{4C}{\rho}\right\rfloor+1\right)B.
\]
Both \cref{eq:movement-finite,eq:movement-lower} also hold with $v$
replaced by $C/L$.
The bounds hold for every ordering of the queries and every matching of
source positions between successive queries.  They also hold when $V$
is measured by ordinary distance on the input line.
\end{theorem}

\begin{outline}
Set $\delta=\rho/4$.  Let $a_j$ be the cumulative maximum-slot movement
before query $j$, so $a_1=0$ and $a_J=C$.  Give query $j$ the label
$\lfloor a_j/\delta\rfloor$.  There are at most
$\lfloor C/\delta\rfloor+1$ labels.  Two queries with the same label
have cumulative movement differing by less than $\delta$; the triangle
inequality therefore keeps each source within distance $\delta$ of its
position at any chosen representative query with that label.

Every centre queried with that label is consequently within $5\rho/4$
of one of the $m$ representative positions.  A circular ball of that
radius contains at most
\[
 \min\left\{N,\left\lfloor\frac{5\rho N}{2L}\right\rfloor+1\right\}
\]
of the equally spaced centres.  The formal proof lifts each covered
centre to a point of the real lattice of spacing $L/N$ inside an interval
of length $5\rho/2$.  Distinct centre indices give distinct lattice
points, whose count is at most the displayed floor bound; the count is
also at most $N$.  The union of the balls contains at most $s$ centres
and hence at most $B$ query types.  Summing over labels proves the integer
bound.  A nonempty sequence forces $B>0$.  Dropping the floor and
rearranging give the two real bounds with $C/L$; $C\le V$ gives the
original summed-movement bounds.

No step used the ordering of the types.  Any chosen matching of the $m$
positions between successive queries defines coherent labelled paths;
apply the same argument to those paths.  Finally, circle distance never
exceeds line distance, so replacing circle variation by line variation
preserves the bound.
\end{outline}

For a tape construction, this theorem charges the movement of the entire
source bank.  The sharper form requires progress even when simultaneous
slot movements are charged only by their maximum.  Relabelling the bank
slots does not remove the bound.  The shared Lean argument allows any
positive partition width $\delta$ and any nonnegative edge charge that
dominates each slot's movement; the ball radius becomes $\rho+\delta$.

\begin{corollary}[Queries spread over a column interval]
\label{cor:movement-interval}
\formal{SourceCovers/IntervalQueryMovement.lean}{SourceCovers.QueryMovement.interval_movement_max,
SourceCovers.QueryMovement.interval_movement}
Let $I_1,\ldots,I_J$ be disjoint, chronologically ordered intervals in a
finite discrete timeline.  Let $m$ labelled source trajectories have
total circle variation $V$ over that entire timeline, including the
gaps between intervals.  Assign a query type $S_j$ to each interval,
and suppose each of the $F$ distinct types occurs at least $r\ge1$
times.  For a width $w\ge0$ and every $c_i\in S_j$, suppose there is a raw point $a_{i,j}$
with
\[
 \operatorname{dist}_L(a_{i,j},c_i)\le w
\]
that is within distance $\rho>0$ of a source position at some time in
$I_j$.  The position and time may depend on $i$ and $j$.  With
$v=V/L$, $\lambda=\rho/L$, and
\[
 s'=\min\{N,\;m(\lfloor3\lambda N+2wN/L\rfloor+1)\},
 \qquad B'=\binom{s'}k,
\]
one has
\begin{align}
 F-\frac{4v}{\lambda r}&\le(1+4v/\lambda)B',
   \label{eq:movement-interval-count}\\
 v&\ge\frac{\lambda}{4}\frac{F-B'}{B'+1/r}.
   \label{eq:movement-interval-lower}
\end{align}
These bounds also hold with $v$ replaced by $C/L$, where
\[
 C=\sum_{t=0}^{T-1}\max_{1\le\ell\le m}
       \operatorname{dist}_L(y_\ell(t+1),y_\ell(t))
\]
for the full timeline $\{0,\ldots,T\}$.  More precisely,
\[
 F-\frac{4C}{\rho r}
 \le\left(\left\lfloor\frac{4C}{\rho}\right\rfloor+1\right)B'.
\]
It suffices that the intervals' internal edge sets are pairwise disjoint;
their endpoints may coincide, and no ordering is needed.  The repetition
bound $r$ may be any positive real number.
\end{corollary}

\begin{outline}
Set $\delta=\rho/4$ and charge each time step by the largest movement of
any source.  Call an interval bad if its internal charge exceeds $\delta$.
Disjoint internal edge sets give at most $C/\delta$ bad intervals.
Every type lost on removing them has all its at least $r$ occurrences
among the deleted intervals.  Counting those disjoint occurrence sets
bounds the number of lost types by $C/(\delta r)$.

Freeze the positions at the start of each remaining interval.  A
required centre is within distance $\rho+\rho/4+w$ of one of these
frozen positions: the raw point is within $w$ of its centre and $\rho$
of a source at a time whose movement from the interval start costs at
most $\delta$.  Apply the cumulative-clock count of \cref{thm:movement}
to the good interval starts, using the full timeline clock and width
$\delta$.  Each label's queried centres lie within $3\rho/2+w$ of
one representative bank, so there are at most $B'$ types per label and
at most $\lfloor C/\delta\rfloor+1$ labels.  The count applies also to
an empty set of good intervals.

Adding the lost types gives the sharper bound.  Dropping the floor and
collecting terms gives both displayed bounds with $C/L$ in place of
$v$; the denominator $B'+1/r$ is positive even when $B'=0$.
Finally $C\le V$ yields the original bounds.  If the original interval
family is empty, the result is immediate.
\end{outline}

Thus queries made at different times during an output column do not
permit us to freeze their positions for free.  Their internal movement
and the possible loss of query types both appear in the bound.
The reusable proof allows independent positive thresholds for internal
interval movement and for grouping interval starts; their sum replaces
the two occurrences of $\rho/4$ in the enlarged radius.

\subsection{A coded-caching calibration}

We use the transpose of the subset placement-delivery array in
\cite[Section IV, Theorem 2]{YanChengTangChen2017}.  Repeating its rows
and columns gives the square matrix below.  We include the elementary
decoding and entropy calculation to specify exactly which interface is
being tested.

Fix integers $N>k\ge2$ and $h\ge1$, and put
\[
 F=\binom Nk,\qquad d=NFh,\qquad L=d^2b,
 \qquad W=db,\qquad a=L/N,
\]
where $b\ge1$ is an integer.  The asymptotic application permits
$b=b(d)$.  Index the rows by $(i,u)\in\{0,\ldots,N-1\}\times
\{0,\ldots,Fh-1\}$, grouping each fixed $i$ into one contiguous bin.
Index columns by $(S,z)$, where $|S|=k$ and $0\le z<Nh$.  A column
decoder may observe all raw rows in the bins $i\in S$, together with
one common helper $G$.  The helper contains the following $b$-bit
parities, for every $(k+1)$-element set $T$ and every $u,z$:
\begin{equation}
 Z_{T,u,z}=\bigoplus_{i\in T}
              X_{(i,u),(T\setminus\{i\},z)}.
 \label{eq:subset-helper}
\end{equation}

\begin{proposition}[Exact helper and gap-area calibration]
\label{prop:subset-calibration}
\notformalized
The specified interface recovers every column on the full input cube.
For a uniform input, the helper is uniform on $g$ bits, where
\begin{equation}
 g=H(G)=L\frac{N-k}{N(k+1)}.
 \label{eq:subset-helper-size}
\end{equation}
Every helper fibre has size $2^{L-g}$.  Every entry with $i\in S$ is
absent from all helper parities and is independent of $G$ and all other
input entries.

For $P_S=\{ia:i\in S\}$ on the raw circle,
\begin{equation}
 \frac1{\binom Nk}\sum_{|S|=k}\gaparea(P_S)
   =\frac{2N-k+1}{2N(k+1)}
   =\frac gL+\frac1{2N}.
 \label{eq:subset-gap-area}
\end{equation}
Consequently, for fixed $k$ and $N\to\infty$, the queried row fraction
$k/N$ tends to zero, whereas both $g/L$ and the mean point-gap area
tend to $1/(k+1)$.
\end{proposition}

\begin{informalproof}
If $i\in S$, the desired entry is read directly.  If $i\notin S$,
take $T=S\cup\{i\}$ in \cref{eq:subset-helper}.  All the other terms
have row-bin index in $S$, so the queried rows and that parity determine
the entry.  Every nonqueried entry belongs to exactly one parity,
specified by $T=S\cup\{i\}$ and its $u,z$.  The parity classes are
disjoint, so their bits are independent and uniform.  Counting them
gives
\[
 g=b\binom N{k+1}(Fh)(Nh)
   =L\frac{N-k}{N(k+1)}.
\]
The parity map is surjective, with all fibres of size $2^{L-g}$.
Queried entries occur in none of these parity classes, proving the
independence assertion.

For the gap formula, choose $S$ uniformly and mark one of its elements
uniformly.  Starting at the marked element, the clockwise integer gaps
$D_1,\ldots,D_k$ are uniform among positive compositions of $N$.
Indeed, each composition and starting point specifies one marked set,
and $N\binom{N-1}{k-1}=k\binom Nk$.  Put $E_i=D_i-1$ and $n=N-k$.
Then $\sum_iE_i=n$.  Stars and bars gives
\[
 \mathbb E\sum_i E_i(E_i-1)
 =\frac{2k\binom{n+k-1}{k+1}}{\binom{n+k-1}{k-1}}
 =\frac{2n(n-1)}{k+1},
\]
with both sides zero when $n<2$.  The numerator follows, for example,
by taking the coefficient of $z^n$ in
$2kz^2(1-z)^{-k-2}$.  Expanding $\sum_i(E_i+1)^2$ now yields
\[
 \mathbb E\sum_iD_i^2=\frac{N(2N-k+1)}{k+1}.
\]
The geometric gaps are $aD_i$, so division by $2L^2$ proves
\cref{eq:subset-gap-area}.  The limits follow from the displayed
formulas.
\end{informalproof}

This example prevents an information-only inference from sparse queried
rows to a deficit exceeding the helper cost.  The exact recovery code
exists, and the normalized gap area approaches the same value as that
cost.  The remaining issue is its geometric realization.

\begin{corollary}[The dimension limit precedes the radius limit]
\label{cor:subset-movement}
\notformalized
Fix $k\ge2$, $m\ge1$ and $h\ge1$, and take the preceding interfaces
with $N\to\infty$.  For each $N$, allow arbitrary column order and
arbitrary coherent source matching.  Let $v_N$ be the total circle
variation divided by $L$.

If every queried bin centre is covered at the column representative
within radius $\rho=\lambda L$, then, for every fixed
$0<\lambda<2/(5m)$,
\begin{equation}
 \liminf_{N\to\infty}v_N
 \ge\frac{\lambda}{4}
      \left[\left(\frac{5m\lambda}{2}\right)^{-k}-1\right].
 \label{eq:subset-centre-limit}
\end{equation}
If instead one required raw point within $W$ of each queried bin centre
is covered within radius $\lambda L$ at some time during its column
interval, then, for every fixed $0<\lambda<1/(3m)$,
\begin{equation}
 \liminf_{N\to\infty}v_N
 \ge\frac{\lambda}{4}
      \bigl[(3m\lambda)^{-k}-1\bigr].
 \label{eq:subset-interval-limit}
\end{equation}
Both right-hand sides tend to infinity as $\lambda\downarrow0$.
Thus no fixed variation constant represents this interface with a fixed
number of source trajectories at every sufficiently small fixed radius
ratio.
\end{corollary}

\begin{informalproof}
For the first assertion, $s/N\to5m\lambda/2<1$ in
\cref{thm:movement}.  Since $k$ is fixed,
\[
 \frac{\binom{s}{k}}{\binom Nk}\longrightarrow
        (5m\lambda/2)^k.
\]
Apply \cref{eq:movement-lower}.  For the second assertion, each type
occurs $r=Nh$ times and $w=W=L/d$.  Hence
$2wN/L=2/(Fh)\to0$, $s'/N\to3m\lambda<1$, and
$\binom{s'}k/\binom Nk\to(3m\lambda)^k$.  Divide the numerator and
denominator of \cref{eq:movement-interval-lower} by $\binom Nk$ to
obtain \cref{eq:subset-interval-limit}.  Each lower bound is a positive
constant times $\lambda^{1-k}$ minus $\lambda/4$, which proves the
last assertion.
\end{informalproof}

The necessary-point premise has a precise information-theoretic use.
Choose, for each queried bin, one bit of an own-column entry in a middle
row.  Its raw coordinate is within $W$ of the bin centre.  That bit is
independent of the specified helper and every other raw bit, so a
decoder on the full cube cannot omit it from its raw dependency set.
Any geometric representation of that dependency set must cover it.
When those raw observations are represented during the column interval,
\cref{cor:subset-movement} applies.

This necessity assertion does not survive every change of helper.
Revealing all entries with $i\in S$ costs exactly $kL/N=o(L)$ extra
bits at fixed $k$ and removes these particular witnesses.  That payment
does not show that the remaining recovery is cheap; it shows why the
movement result is about the stated interface.  A native source-bank
argument may reveal additional threshold or residual information, so
its required queries must be established after those exact disclosures.
The radius ratio is also fixed before the dimension limit; the
conclusion does not interchange these limits by setting $\lambda=1/N$.
